% !PRINTED 355
\weberpage{355}

%<<CONT
eine durch $\mathfrak{m}$ teilbare ganze rationale Zahl ist, so muß sie durch $a$ teilbar sein; setzen wir sie $= ac$, so ist $c$ eine ganze rationale Zahl und
\weq{6}{D = b^2 - 4ac,}
\weq{7}{b^2 \equiv D \pmod{4a},}
und wenn die Kongruenz (7) erfüllt ist, so genügt auch jede mit $b$ nach dem Modul $2a$ kongruente Zahl $b'$ derselben Kongruenz.

Wir schließen hieraus:

5. \gesperrt{Satz. Jedes primäre zu $Q$ teilerfremde Ideal $\mathfrak{m}$, dessen Norm $= a$ ist, liefert uns durch (4) eine und nur eine Wurzel $b$ der Kongruenz (7), wenn als Wurzeln nicht einzelne Zahlen, sondern Zahlklassen nach dem Modul $2a$ betrachtet werden.}

Die beiden Zahlen
\weq{8}{a, \quad \frac{b + \sqrt{D}}{2}}
können, wenn $a$ relativ prim zu $Q$ ist, nicht beide durch eine und dieselbe rationale Primzahl teilbar sein; denn sonst müßte
\[ \frac{b + \sqrt{D}}{2} - \frac{b - \sqrt{D}}{2} = \sqrt{D} \]
durch $p$ teilbar sein; also wäre $D$ durch $p^2$ teilbar und $p$ müßte ein Teiler von $Q$ sein, gegen die Voraussetzung.

Für ein ungerades $p$ ist dies evident, für $p = 2$ aber würde folgen:
\[ b \equiv 0 \pmod{2}, \quad b^2 - D \equiv 0 \pmod{16}, \]
und mithin
\[ D \equiv 0, 4 \pmod{16}. \]
Es wäre also $D/4$ noch Diskriminante und 2 müßte in $Q$ aufgehen. Damit ist bewiesen:

6. Satz. \gesperrt{Ist $a$ eine zu $Q$ teilerfremde natürliche Zahl und $b$ eine Wurzel von (7), so erhält man ein bestimmtes primäres Ideal $\mathfrak{m}$ als größten gemeinschaftlichen Teiler der beiden Zahlen (8).}

Jetzt beweisen wir:

7. \gesperrt{Satz. Jede durch das primäre Ideal $\mathfrak{m}$ teilbare, zu $Q$ teilerfremde Zahl $\mu$ der Ordnung $[Q]$ ist in der Form}
\weq{9}{\mu = ax + \frac{b + \sqrt{D}}{2} y}
\gesperrt{einmal und nur einmal darstellbar, wenn $x, y$ ganze rationale Zahlen sind. Werden $x, y$ als Variable be-}
%CONT>>
